\chapter[Poprawność rachunku pierwszego rzędu]{Poprawność rachunku dowodowego i konsekwencja logiczna}
\label{ch:soundness-fol}

\section{Dwa rodzaje wynikania}
Zapis $\Gamma\vdash_{\mathrm{ND}}\varphi$ z
rozdziału~\ref{ch:natural-deduction} oznacza istnienie
skończonego, formalnego dowodu. Zapis $\Gamma\models\varphi$
z rozdziału~\ref{ch:semantics-fol} mówi natomiast o
prawdziwości formuły we wszystkich interpretacjach spełniających założenia.
Są to pojęcia zbudowane odmiennymi metodami.
Pierwsze dotyczy skończonych napisów i reguł, drugie
— wszystkich struktur i ich wartościowań.

\begin{definicja}[Konsekwencja semantyczna dla formuł]
\label{def:soundness-entailment-formulas}
Niech $\mathcal L$ będzie językiem pierwszego rzędu,
$\Gamma$ zbiorem formuł języka $\mathcal L$,
a $\varphi$ formułą tego samego języka.
Piszemy $\Gamma\models\varphi$, jeżeli dla każdej
struktury $\mathfrak A$ języka $\mathcal L$
o niepustym uniwersum $A$ i każdego wartościowania $s$ zmiennych języka
$\mathcal L$ w zbiorze $A$ zachodzi implikacja
\[
\bigl(\forall\gamma\in\Gamma\ \mathfrak A,s\models\gamma\bigr)
\quad\Longrightarrow\quad
\mathfrak A,s\models\varphi.
\]
Gdy $\Gamma=\varnothing$, piszemy $\models\varphi$.
\end{definicja}
Jeżeli $\Gamma$ składa się ze zdań, a $\varphi$ jest zdaniem,
to dzięki Wnioskowi~\ref{cor:sem-sentence}
powyższa definicja zgadza się z
Definicją~\ref{def:sem-entailment}.
Nie można jednak pominąć wartościowania, gdy w formułach występują
zmienne wolne. Przykładowo $P(x)$ nie wynika semantycznie
z samego istnienia elementu spełniającego $P$.

\begin{przyklad}[Różnica między pojedynczym modelem a wszystkimi modelami]
W strukturze $(\N,<)$ zdanie $\forall x\,\exists y\,(x<y)$
jest prawdziwe. Nie oznacza to, że jest prawdziwe
w każdej strukturze z binarnym symbolem relacyjnym $<$.
Jeśli za uniwersum przyjmiemy $\{a\}$, a relację $<$
za pustą, to zdanie jest fałszywe. Dla wynikania
semantycznego potrzebna jest kwantyfikacja po \emph{wszystkich}
strukturach rozważanego języka.
\end{przyklad}


\begin{przyklad}[Semantyczna wersja prostego kroku dowodowego]
Sprawdzając semantyczną poprawność reguły, nie twierdzimy,
że jej przesłanki są prawdziwe w każdej strukturze.
Pytamy jedynie: \emph{jeżeli} w ustalonej strukturze
przesłanki są spełnione, czy musi być spełniony wniosek?
To pozwala zrozumieć, dlaczego nad kreską zapisujemy
kilka wymaganych przesłanek, nie utożsamiając ich
z dodatkowymi aksjomatami całej teorii.

Niech $P$ i $Q$ oznaczają dowolne formuły tego samego języka.
Mamy $\{P\land Q\}\models P$: w dowolnej strukturze i przy
dowolnym wartościowaniu, w którym $P\land Q$ jest prawdziwa,
prawdziwe jest także $P$. To jest fakt o \emph{znaczeniu}
spójnika $\land$. W rozdziale~\ref{ch:natural-deduction}
ten sam krok zapisaliśmy formalną regułą $(\land E_1)$.
Twierdzenie o poprawności pokaże, że tak wygląda każdy
legalny krok systemu dowodowego: nie może zamienić
prawdziwych przesłanek na fałszywy wniosek.
\end{przyklad}

\section{Trzy pomocnicze fakty semantyczne}
Dwa pierwsze fakty udowodniliśmy już w rozdziale 10,
ale przytoczymy je w dokładnej postaci potrzebnej poniżej.

\begin{lemat}[Zmiana jednej nieistotnej zmiennej]
\label{lem:soundness-fresh}
Niech $\mathcal L$ będzie językiem pierwszego rzędu,
$\mathfrak A$ dowolną strukturą tego języka o uniwersum $A$,
$\varphi$ formułą, a $y$ zmienną niewystępującą wolno
w $\varphi$. Dla każdego wartościowania $s$ i $a\in A$ zachodzi
\[
\mathfrak A,s\models\varphi
\quad\Longleftrightarrow\quad
\mathfrak A,s[y\mapsto a]\models\varphi.
\]
\end{lemat}
\begin{proof}
Wartościowania $s$ i $s[y\mapsto a]$ różnią się
co najwyżej na zmiennej $y$. Ponieważ $y$ nie jest
wolną zmienną $\varphi$, zgadzają się na
$\operatorname{FV}(\varphi)$. Teza wynika z
Twierdzenia~\ref{thm:sem-coincidence}.
\end{proof}

\begin{lemat}[Interpretacja podstawienia]\label{lem:soundness-subst}
Niech $\mathcal L$ będzie językiem pierwszego rzędu,
$\mathfrak A$ jego strukturą o uniwersum $A$,
$\varphi$ formułą, $x$ zmienną, $t$ termem dopuszczalnym
do podstawienia za $x$ w $\varphi$, a $s$ wartościowaniem.
Wówczas
\[
\mathfrak A,s\models\varphi[x:=t]
\quad\Longleftrightarrow\quad
\mathfrak A,s[x\mapsto t^{\mathfrak A}[s]]\models\varphi.
\]
\end{lemat}
\begin{proof}
To Twierdzenie~\ref{thm:sem-substitution},
którego dowód indukcyjny obejmował wszystkie rodzaje formuł,
w szczególności przypadki obu kwantyfikatorów.
\end{proof}

\begin{lemat}[Zastępowanie równych termów]
\label{lem:soundness-equality}
Niech $\mathcal L$ będzie językiem pierwszego rzędu
z równością, $\mathfrak A$ strukturą języka $\mathcal L$
o niepustym uniwersum $A$, $s$ wartościowaniem,
$x$ zmienną, $t,u$ termami języka $\mathcal L$,
a $\theta$ formułą tego języka.
Załóżmy, że $t$ i $u$ są dopuszczalne do podstawienia
za zmienną $x$ w $\theta$.
Jeżeli zachodzą $\mathfrak A,s\models t=u$ oraz
$\mathfrak A,s\models\theta[x:=t]$, to
$\mathfrak A,s\models\theta[x:=u]$.
\end{lemat}
\begin{proof}
Z prawdziwości atomu $t=u$ otrzymujemy
$t^{\mathfrak A}[s]=u^{\mathfrak A}[s]=a$
dla pewnego $a\in A$.
Lemat~\ref{lem:soundness-subst} przekształca przesłankę
$\mathfrak A,s\models\theta[x:=t]$ w
$\mathfrak A,s[x\mapsto a]\models\theta$.
Ponieważ $u^{\mathfrak A}[s]=a$, ponowne zastosowanie
tego lematu daje
$\mathfrak A,s\models\theta[x:=u]$.
\end{proof}

\section{Twierdzenie o poprawności}
\begin{twierdzenie}[Poprawność dedukcji naturalnej]
\label{thm:soundness}
Niech $\mathcal L$ będzie dowolnym językiem pierwszego rzędu
z równością, $\Gamma$ zbiorem formuł tego języka,
a $\varphi$ formułą. Jeśli
\[
\Gamma\vdash_{\mathrm{ND}}\varphi,
\]
to
\[
\Gamma\models\varphi.
\]
\end{twierdzenie}
\begin{proof}
Ustalmy dowolną strukturę $\mathfrak A$ języka $\mathcal L$
i dowolne wartościowanie $s$ spełniające wszystkie
formuły z $\Gamma$. Pokażemy, że spełniona jest także $\varphi$.
Dowód prowadzimy indukcją po strukturze skończonego drzewa dedukcji.
Hipoteza indukcyjna mówi dokładniej, że każda poddedukcja
jest semantycznie poprawna względem swoich \emph{otwartych}
założeń. Rozładowane założenia rozpatrujemy tylko na czas
odpowiedniego kroku.

\smallskip
\emph{1. Liście i reguły bez przesłanek.}
Jeśli liść jest otwartym założeniem $\gamma$
danej poddedukcji, to każde wartościowanie
spełniające wszystkie jej otwarte założenia spełnia też
$\gamma$. Dotyczy to również założeń tymczasowych,
które będą rozładowane w dalszych krokach.
Dla całego dowodu jego otwarte założenia należą do $\Gamma$.
Reguła $=I$ kończąca się na $t=t$ jest poprawna,
ponieważ $t^{\mathfrak A}[s]=t^{\mathfrak A}[s]$.

\smallskip
\emph{2. Koniunkcja i implikacja.}
Jeżeli obie przesłanki reguły $\land I$ są spełnione,
to prawdziwa jest ich koniunkcja; jeśli spełniona jest
koniunkcja, każdy jej człon jest spełniony, co uzasadnia
$\land E_1$ i $\land E_2$.
Dla $\to E$ spełnianie $\alpha\to\beta$ oraz $\alpha$
wymusza spełnianie $\beta$ zgodnie z tablicą implikacji.
Dla $\to I$ mamy poddedukcję $\beta$ z dodatkowym,
tymczasowym założeniem $\alpha$. Rozważmy dowolną
strukturę i wartościowanie spełniające pozostałe otwarte
założenia. Jeśli $\alpha$ jest fałszywa, implikacja
$\alpha\to\beta$ jest prawdziwa niezależnie od $\beta$.
Jeśli $\alpha$ jest prawdziwa, hipoteza indukcyjna
zastosowana do poddedukcji daje $\beta$.
W obu przypadkach implikacja jest prawdziwa.

\smallskip
\emph{3. Alternatywa, równoważność i negacja.}
Reguły $\lor I_1$ i $\lor I_2$ są poprawne, bo
spełnianie jednego członu wystarcza dla alternatywy.
W $\lor E$ prawdziwa alternatywa oznacza prawdziwość
co najmniej jednego członu. W pierwszym przypadku
hipoteza indukcyjna dla pierwszej poddedukcji daje $\theta$;
w drugim przypadku daje ją druga poddedukcja.
Reguły równoważności są poprawne, ponieważ
$\alpha\leftrightarrow\beta$ jest prawdziwa dokładnie wtedy,
gdy prawdziwe są obie implikacje wzajemne.

Jeśli prawdziwe byłyby zarazem $\alpha$ i $\neg\alpha$,
otrzymalibyśmy sprzeczność z definicją negacji.
Dlatego reguła $\neg E$ może prowadzić do $\bot$
jedynie z przesłanek, których nie można jednocześnie spełnić.
Jeśli z dodatkowego założenia $\alpha$ otrzymano $\bot$,
to przy spełnieniu pozostałych założeń $\alpha$ nie może być
prawdziwa: inaczej otrzymalibyśmy spełnianie $\bot$.
Zatem $\neg\alpha$ jest prawdziwa, co uzasadnia $\neg I$.
Reguła $\bot E$ jest poprawna, ponieważ $\bot$ nie jest
spełniona przy żadnym wartościowaniu; zatem nie istnieje
wartościowanie, dla którego jej przesłanka byłaby prawdziwa,
a wniosek fałszywy. Wreszcie reguła klasycznego dowodu
nie wprost $\mathrm{RAA}$ jest poprawna: jeśli założenie
$\neg\alpha$ prowadziłoby do sprzeczności, nie może być
prawdziwe, a klasyczna semantyka dwuwartościowa
zapewnia prawdziwość $\alpha$.

\smallskip
\emph{4. Eliminacja kwantyfikatora ogólnego.}
Przypuśćmy, że $\mathfrak A,s\models\forall x\,\alpha$.
Niech $a=t^{\mathfrak A}[s]$. Definicja spełniania
kwantyfikatora daje
$\mathfrak A,s[x\mapsto a]\models\alpha$.
Lemat~\ref{lem:soundness-subst} zamienia ten warunek
na $\mathfrak A,s\models\alpha[x:=t]$.
To dokładnie wniosek reguły $\forall E$.

\smallskip
\emph{5. Wprowadzanie kwantyfikatora ogólnego.}
Niech poddedukcja z otwartymi założeniami $\Delta$
kończy się na $\alpha$, a zmienna $x$ nie występuje
wolno w żadnej formule z $\Delta$.
Załóżmy, że $\mathfrak A,s$ spełnia $\Delta$.
Dla dowolnego $a\in A$ Lemat~\ref{lem:soundness-fresh}
zastosowany oddzielnie do każdego elementu $\Delta$
pokazuje, że $\mathfrak A,s[x\mapsto a]$ również spełnia
wszystkie otwarte założenia. Hipoteza indukcyjna daje więc
$\mathfrak A,s[x\mapsto a]\models\alpha$.
Ponieważ $a$ było dowolne, definicja $\forall$ daje
$\mathfrak A,s\models\forall x\,\alpha$.
W tym właśnie miejscu konieczny jest warunek o niewystępowaniu
$x$ w otwartych założeniach.

\smallskip
\emph{6. Wprowadzanie kwantyfikatora egzystencjalnego.}
Jeśli $\mathfrak A,s\models\alpha[x:=t]$,
to dla $a=t^{\mathfrak A}[s]$ Lemat~\ref{lem:soundness-subst}
daje $\mathfrak A,s[x\mapsto a]\models\alpha$.
Wartość $a$ jest świadkiem spełniania $\exists x\,\alpha$.

\smallskip
\emph{7. Eliminacja kwantyfikatora egzystencjalnego.}
Załóżmy, że ostatnia reguła ma przesłanki
$\exists x\,\alpha$ oraz poddedukcję $\theta$ z dodatkowego
założenia $\alpha[x:=y]$, gdzie $y$ jest świeżą zmienną
nieobecną w $\alpha$ i we wniosku $\theta$ oraz niewystępującą
wolno w żadnym otwartym założeniu obu gałęzi poza
rozładowywanym $\alpha[x:=y]$. Przez $\Delta$ oznaczmy
sumę tych pozostałych otwartych założeń; w szczególności
obejmuje ona założenia potrzebne do wyprowadzenia
przesłanki głównej $\exists x\,\alpha$.
Niech $\mathfrak A,s$ spełnia $\Delta$.
Hipoteza indukcyjna zastosowana do dowodu przesłanki głównej
zapewnia $\mathfrak A,s\models\exists x\,\alpha$.
Z tego wynika istnienie
$a\in A$ takiego, że
$\mathfrak A,s[x\mapsto a]\models\alpha$.
Połóżmy $r=s[y\mapsto a]$.
Ponieważ $y$ nie występuje wolno w $\alpha$,
wartościowania $r[x\mapsto a]$ i $s[x\mapsto a]$
zgadzają się na wszystkich jej zmiennych wolnych.
Stąd na mocy Twierdzenia~\ref{thm:sem-coincidence}
mamy $\mathfrak A,r[x\mapsto a]\models\alpha$.
Z Lematu~\ref{lem:soundness-subst} dla termu $y$ wynika
$\mathfrak A,r\models\alpha[x:=y]$.
Świeżość $y$ oraz Lemat~\ref{lem:soundness-fresh}
zapewniają też, że $\mathfrak A,r$ spełnia $\Delta$.
Hipoteza indukcyjna dla poddedukcji daje zatem
$\mathfrak A,r\models\theta$. Zmienna $y$ nie występuje
wolno w $\theta$, więc ponowne użycie lematu o świeżej
zmiennej daje $\mathfrak A,s\models\theta$.
Reguła $\exists E$ jest więc poprawna.

\smallskip
\emph{8. Eliminacja równości.}
Jest to dokładnie treść Lematu~\ref{lem:soundness-equality}.

Wyczerpaliśmy wszystkie reguły z
Definicji~\ref{def:nd-rules-conj-imp},
Definicji~\ref{def:nd-rules-prop},
Definicji~\ref{def:nd-rules-quant} i
Definicji~\ref{def:nd-equality}.
Indukcja po drzewie dowodu jest zakończona.
Dla korzenia drzewa otwarte założenia zawierają się w $\Gamma$,
więc przyjęte na początku $\mathfrak A,s$ spełnia także wniosek
$\varphi$. Ponieważ $\mathfrak A$ i $s$ były dowolne,
wynika $\Gamma\models\varphi$.
\end{proof}


\section{Jak poprawność działa przy pominiętych krokach?}
Dla czytelnika przyzwyczajonego do zapisu $\vdots$
warto oddzielić trzy poziomy rozumowania.

\begin{przyklad}[Pełne uzasadnienie krótkiego drzewa]
Rozważmy udowodnioną w rozdziale 11 implikację
$(P\land Q)\to(Q\land P)$. Jeżeli w dowolnej strukturze
przy danym wartościowaniu prawdziwa jest formuła $P\land Q$,
to na mocy definicji koniunkcji prawdziwe są osobno $P$ i $Q$.
Wobec tego prawdziwa jest $Q\land P$. Gdy $P\land Q$
nie jest prawdziwa, implikacja jest prawdziwa ze względu na
swoją tabelę prawdy. Oznacza to, że cała implikacja jest
prawdziwa przy każdym wartościowaniu i w każdej strukturze.

W drzewie formalnym pomiędzy $[P\land Q]^1$ i $Q\land P$
można było użyć skrótu $\vdots$, ale jego poprawność
weryfikujemy przez sprawdzenie trzech konkretnych zastosowań
reguł: $(\land E_2)$, $(\land E_1)$ i $(\land I)$.
Twierdzenie~\ref{thm:soundness} mówi, że taka weryfikacja
może być przeprowadzona dla \emph{każdego} skończonego drzewa.
\end{przyklad}

\begin{uwaga}[Poprawność nie mówi, że założenia są prawdziwe]
Jeżeli $\Gamma\vdash_{\mathrm{ND}}\varphi$, poprawność orzeka:
\emph{w każdej interpretacji, w której spełnione są wszystkie
formuły z $\Gamma$, spełniona jest także $\varphi$}.
Jeżeli $\Gamma$ ma jakieś fałszywe założenie w danej
interpretacji, twierdzenie nie daje w tej interpretacji
informacji o wartości $\varphi$. Otwieranie i rozładowywanie
założeń regułą $(\to I)$ ma więc sens \emph{warunkowy},
nie zmienia wartości logicznych formuł.
\end{uwaga}

\section{Konsekwencje poprawności}
\begin{wniosek}[Brak formalnego dowodu zdania fałszywego w pewnym modelu]
\label{cor:soundness-countermodel}
Niech $\mathcal L$ będzie językiem pierwszego rzędu,
$\Gamma$ zbiorem jego formuł, a $\varphi$ formułą
tego języka. Jeśli istnieją struktura $\mathfrak A$
języka $\mathcal L$ i wartościowanie $s$
spełniające wszystkie formuły z $\Gamma$, ale niespełniające
$\varphi$, to $\Gamma\nvdash_{\mathrm{ND}}\varphi$.
\end{wniosek}
\begin{proof}
Gdyby $\Gamma\vdash_{\mathrm{ND}}\varphi$,
Twierdzenie~\ref{thm:soundness} dawałoby
$\Gamma\models\varphi$. To przeczyłoby istnieniu
opisanego kontrmodelu.
\end{proof}
\begin{przyklad}[Dlaczego $\forall x\,P(x)$ nie wynika z $P(y)$]
Niech uniwersum struktury wynosi $\{0,1\}$,
a $P^{\mathfrak A}=\{0\}$. Przy wartościowaniu
$s(y)=0$ spełniona jest formuła $P(y)$,
ale nie jest spełnione $\forall x\,P(x)$.
Z Wniosku~\ref{cor:soundness-countermodel} otrzymujemy
$\{P(y)\}\nvdash_{\mathrm{ND}}\forall x\,P(x)$.
To samo można było zauważyć bezpośrednio na poziomie
reguł: warunek dla $\forall I$ jest naruszony, jeśli
uogólniamy zmienną wolną w nieusuniętym założeniu.
\end{przyklad}

\begin{wniosek}[Model wyklucza dowód sprzeczności]
\label{cor:soundness-consistency}
Niech $\mathcal L$ będzie językiem pierwszego rzędu,
a $\Gamma$ zbiorem formuł tego języka.
Jeśli istnieją struktura $\mathfrak A$ języka $\mathcal L$
i wartościowanie $s$ spełniające wszystkie
formuły z $\Gamma$, to
$\Gamma\nvdash_{\mathrm{ND}}\bot$.
\end{wniosek}
\begin{proof}
Formuła $\bot$ jest fałszywa w każdej strukturze.
Model $\Gamma$ jest więc kontrmodelem dla
$\Gamma\models\bot$. Z poprzedniego wniosku otrzymujemy
$\Gamma\nvdash_{\mathrm{ND}}\bot$.
\end{proof}

\begin{uwaga}[Ostrożnie z odwracaniem implikacji]
Udowodniliśmy: istnienie modelu daje brak formalnego dowodu
sprzeczności. Odwrotność — że każdy niesprzeczny zbiór zdań
ma model — nie wynika z samej poprawności.
Wymaga twierdzenia o pełności, a dla nieskończonych zbiorów zdań
stanowi jego ważną postać.
\end{uwaga}

\section{Twierdzenia o dedukcji: syntaktyczne i semantyczne}
\begin{twierdzenie}[Semantyczne twierdzenie o dedukcji]
\label{thm:soundness-semantic-deduction}
Niech $\mathcal L$ będzie językiem pierwszego rzędu,
$\Gamma$ zbiorem formuł tego języka, a $\alpha,\beta$
formułami języka $\mathcal L$. Wtedy
\[
\Gamma\cup\{\alpha\}\models\beta
\quad\Longleftrightarrow\quad
\Gamma\models\alpha\to\beta.
\]
\end{twierdzenie}
\begin{proof}
Załóżmy, że $\Gamma\cup\{\alpha\}\models\beta$.
Ustalmy strukturę $\mathfrak A$ i wartościowanie $s$
spełniające $\Gamma$. Jeśli nie spełniają $\alpha$,
implikacja $\alpha\to\beta$ jest prawdziwa.
Jeśli spełniają $\alpha$, to spełniają wszystkie założenia
z $\Gamma\cup\{\alpha\}$, a więc także $\beta$.
Implikacja jest znów prawdziwa. Zatem
$\Gamma\models\alpha\to\beta$.

W drugą stronę załóżmy
$\Gamma\models\alpha\to\beta$.
Niech $\mathfrak A,s$ spełnia wszystkie formuły z
$\Gamma\cup\{\alpha\}$. Spełnia więc $\Gamma$
i $\alpha$. Z założenia wynika spełnianie implikacji
$\alpha\to\beta$. Ponieważ jej poprzednik jest prawdziwy,
na mocy tablicy implikacji prawdziwy jest następnik $\beta$.
To dowodzi $\Gamma\cup\{\alpha\}\models\beta$.
\end{proof}

Zestawmy to z Twierdzeniem~\ref{thm:nd-deduction},
które dotyczyło dowodliwości syntaktycznej. Obie równoważności
mają podobną postać, ale zostały udowodnione na podstawie
zupełnie różnych definicji: jedna za pomocą reguł drzewa dowodu,
druga przez rozważenie dowolnej struktury i wartościowania.

\section{Skończoność formalnych dowodów a zwartość semantyczna}
\begin{twierdzenie}[Syntaktyczna zwartość wyprowadzalności]
\label{thm:soundness-syntactic-finitary}
Niech $\mathcal L$ będzie językiem pierwszego rzędu,
$\Gamma$ zbiorem formuł tego języka, a $\varphi$
formułą języka $\mathcal L$.
Jeżeli $\Gamma\vdash_{\mathrm{ND}}\varphi$, to istnieje
skończony zbiór $\Gamma_0\subseteq\Gamma$ taki, że
$\Gamma_0\vdash_{\mathrm{ND}}\varphi$.
\end{twierdzenie}
\begin{proof}
Twierdzenie~\ref{thm:nd-finite-support} dowiodło tej tezy
bez odwoływania się do semantyki. Powodem jest skończoność
liczby liści drzewa formalnego dowodu.
\end{proof}

\begin{uwaga}[Czego \emph{nie} wolno z tego wywnioskować]
Z samej powyższej własności i poprawności nie możemy jeszcze
stwierdzić, że jeśli $\Gamma\models\varphi$, to istnieje
skończony $\Gamma_0\subseteq\Gamma$ z
$\Gamma_0\models\varphi$.
Do takiego przejścia trzeba mieć także twierdzenie,
że konsekwencja semantyczna jest zawsze dowodliwa
(pełność pierwszego rzędu).
Nie należy więc utożsamiać \emph{syntaktycznej skończoności
dowodów} z \emph{semantycznym twierdzeniem o zwartości}.
\end{uwaga}

\section{Co dalej?}
Dysponujemy teraz jednym kierunkiem pomostu między składnią
a semantyką:
\[
\Gamma\vdash_{\mathrm{ND}}\varphi
\quad\Longrightarrow\quad
\Gamma\models\varphi.
\]
Przed nami kierunek odwrotny. Jego dowód wymaga skonstruowania
modelu dla każdego odpowiednio niesprzecznego zbioru zdań.
W późniejszych rozdziałach wprowadzimy rozszerzenia teorii
o stałe świadczące o istnieniu, zbudujemy model termowy
i przeprowadzimy szczegółowy dowód twierdzenia Gödla o pełności.
Dopiero potem wyprowadzimy semantyczne twierdzenie o zwartości
i dalej rozwiniemy teorię modeli.